\documentclass[11pt]{article}

\usepackage{series}   % provided by the user
\usepackage{amsmath,amssymb,amsthm,amsfonts}
\usepackage{mathrsfs}
\usepackage{hyperref}
\usepackage{geometry}
\geometry{margin=1in}

\numberwithin{equation}{section}


% operators
\DeclareMathOperator{\Dom}{Dom}
\DeclareMathOperator{\Spec}{Spec}

\begin{document}

\begin{center}
{\Large\bf Twistor Encodings as High-Coherence Limits of Modal Triplet Theory}

\vspace{0.4cm}

{\large A Coherent-Sector Encoding, Controlled Truncation, and Kernel Transition Criterion}

\vspace{0.6cm}

{\normalsize
Peter Nero
}
\end{center}

\vspace{0.6cm}

\begin{abstract}
Penrose twistor theory provides a powerful holomorphic encoding of massless,
self-dual, and conformally invariant sectors of four-dimensional field theories,
yet has long resisted interpretation as a fundamental description of nature.
In this paper we show that twistor theory arises naturally as an \emph{effective
encoding} of a distinguished high-coherence regime of Modal Triplet Theory (MTT).
Using the coherent-sector projection, spectral-gap control, and controlled
truncation framework of MTT, we define a precise \emph{twistor corner} in which
twistor methods are valid, explain their efficiency and limitations, and derive
a sharp criterion for their breakdown.

Technically, we identify the twistor corner as the regime in which noncoherent
modes are suppressed, mass-generation effects are negligible, scale variation is
slow, and the coherent gauge or gravitational sector is near self-dual. We prove
that in this regime the coherent dynamics admits a Schur--Feshbach reduced
generator whose correction term is suppressed by the inverse spectral gap.
Twistor encoding is valid precisely while this controlled truncation correction
remains below an admissibility threshold. When the gap collapses or coherent--noncoherent
mixing grows beyond threshold, the induced evolution on coherent data becomes
kernel-valued, signaling a genuine MTT selection event and forcing re-projection
into a new admissible basin.

This work situates twistor theory within a broader theory of coherence, selection,
and irreversibility, explaining both its remarkable successes and its intrinsic
domain of validity.
\end{abstract}

\tableofcontents

\section{Introduction}

Twistor theory, introduced by Penrose \cite{Penrose1967,Penrose1972}, replaces
spacetime fields by holomorphic data on a complex geometric space encoding null
directions and conformal structure. It has achieved striking successes: the
Penrose transform for massless fields, the Ward correspondence for self-dual
Yang--Mills theory, the nonlinear graviton for self-dual gravity, and the modern
twistor-amplitude revolution initiated by Witten \cite{Witten2004} and developed
further in \cite{CachazoSvrcekWitten2004,MasonSkinner2009}.

At the same time, twistor theory has well-known limitations. It is naturally
adapted to massless and self-dual sectors, struggles with mass generation,
confinement, and generic curved backgrounds, and does not by itself provide a
framework for probability, measurement, or irreversibility. These facts have
often been viewed as obstacles to twistor theory as a fundamental description.

Modal Triplet Theory (MTT) takes a different starting point. Rather than proposing
a new geometric container, MTT introduces a \emph{selection principle} based on
coherent projection, spectral gaps, and controlled truncation. Observable physics
arises as the coherent sector of a higher-dimensional modal dynamics, and
effective descriptions are valid only within admissible slabs where coherence
capacity remains positive.

The central claim of this paper is that twistor theory is neither accidental nor
fundamental. Instead:

\begin{quote}
\emph{Twistor theory is the optimal holomorphic encoding of the maximally coherent,
null-dominated, self-dual corner of the MTT coherent sector.}
\end{quote}

MTT explains why this corner exists, why twistor methods are so effective there,
and why they must fail beyond it.

\subsection*{What this paper does and does not do}

\begin{itemize}
\item We do \emph{not} propose a new twistor theory.
\item We do \emph{not} claim twistor space is fundamental or replaces the MTT modal
space.
\item We \emph{do} provide a precise regime (the twistor corner) in which twistor
encoding is valid.
\item We \emph{do} derive a sharp breakdown criterion tied to MTT admissibility and
controlled truncation.
\end{itemize}

\section{Modal Triplet Theory: Minimal Background}

We summarize only the components of MTT required for this paper. For full
development see the MTT corpus \cite{MTTFoundation,MTTFixedPoints,MTTTruncation,MTTCoherence}.

\subsection{Modal geometry and coherent projection}

MTT is defined on a modal configuration space
\begin{equation}
M_{10} = Y^4 \times B_1 \times B_2 \times B_3 ,
\end{equation}
where $Y^4$ is the emergent spacetime base and $B_i$ are compact internal bundles
equipped with Laplace-type operators $\Delta_{B_i}$. The $\Delta_{B_i}$ commute,
allowing definition of a joint coherent projector
\begin{equation}
\Pi_{\mathrm{coh}} = \Pi_{B_1}\Pi_{B_2}\Pi_{B_3},
\end{equation}
where $\Pi_{B_i}$ is the Riesz projector onto $\ker \Delta_{B_i}$.

Observable physics arises from the coherent sector
\[
H_P := \Ran \Pi_{\mathrm{coh}},
\]
while $H_Q := \Ran (I-\Pi_{\mathrm{coh}})$ contains noncoherent modes.

\subsection{Admissible slabs and controlled truncation}

An \emph{admissible slab} $\mathcal S \subset Y^4$ is a spacetime region on which:

\begin{itemize}
\item the internal Laplacians have a uniform spectral gap $\lambda_Q>0$,
\item the projector $\Pi_{\mathrm{coh}}$ is bounded on the relevant Sobolev scales,
\item truncation errors from discarding $H_Q$ are controlled.
\end{itemize}

On such slabs, effective four-dimensional quantum field theory and general
relativity emerge as controlled approximations. When admissibility fails, no
single-valued effective description exists.

\section{Definition of the Twistor Corner}

We now define the regime in which twistor theory becomes a natural encoding of the
MTT coherent sector.

\begin{definition}[Twistor Corner]
Let $\mathcal S$ be an admissible slab. We say that $\mathcal S$ lies in the
\emph{twistor corner} of MTT if the following dimensionless parameters satisfy
\[
(\varepsilon_1,\varepsilon_2,\varepsilon_3,\varepsilon_4) \ll 1:
\]
\begin{enumerate}
\item (\emph{Self-duality}) $\varepsilon_1 := \|F^-\|/\|F\|$, where $F^\pm$ are the
(anti-)self-dual components of the coherent gauge curvature.
\item (\emph{Mass suppression}) $\varepsilon_2 := m/E$, where $m$ denotes masses
generated by overlap/Higgs structure and $E$ is the characteristic energy scale.
\item (\emph{Slow scale variation}) $\varepsilon_3 := \|\nabla C_{\mathrm{MTT}}\|/
C_{\mathrm{MTT}}$, where $C_{\mathrm{MTT}}$ is the coherence-capacity bookkeeping
field.
\item (\emph{Leakage}) $\varepsilon_4 := \|(I-\Pi_{\mathrm{coh}})\Psi\|/\|\Psi\|$.
\end{enumerate}
\end{definition}

\begin{remark}
The parameters $\varepsilon_i$ are dimensionless control quantities whose smallness
is understood in operator-norm or Sobolev-norm estimates appropriate to the slab
$\mathcal S$. No claim of uniform smallness across all scales or all solutions is
made; rather, $\varepsilon_i \ll 1$ denotes membership in an admissible neighborhood
of the twistor corner on $\mathcal S$.
\end{remark}

\begin{remark}
Intuitively, the twistor corner is the regime in which the coherent dynamics is
null-dominated, approximately conformal, and near an integrable self-dual
subsector, with negligible excitation of noncoherent modes.
\end{remark}

In the remainder of the paper we show that:
\begin{itemize}
\item in the twistor corner, twistor theory provides a faithful and efficient
encoding of the coherent EFT;
\item deviations from this corner correspond to controlled holomorphicity-breaking
deformations;
\item loss of admissibility corresponds to a sharp transition to kernel-valued
evolution.
\end{itemize}

\bigskip

% References will be added in the final section
%%%%%%%%%%%%%%%%%%%%%%%%%%%%%%%%%%%%%%%%%%%%%%%%%%%%%%%%%%%%%%%%%%%%%%%%%%%%%%
% Part II: Generator Decomposition and Controlled Truncation
%%%%%%%%%%%%%%%%%%%%%%%%%%%%%%%%%%%%%%%%%%%%%%%%%%%%%%%%%%%%%%%%%%%%%%%%%%%%%%

\section{Generator Decomposition and Projected Dynamics}

We now place the twistor corner definition into the analytic framework used
throughout the MTT corpus: semigroup evolution, coherent projection, and
controlled truncation.

\subsection{Unprojected and projected evolution}

Let $H$ denote the slab-local state space (a Hilbert or Banach space depending on
the model), and let
\[
\Phi_t = e^{tL}, \qquad t \ge 0,
\]
be the unprojected evolution on $H$ generated by a closed operator
\[
L : \Dom(L) \subset H \to H .
\]

Define the coherent and noncoherent subspaces
\[
H_P := \Ran \Pi_{\mathrm{coh}}, \qquad
H_Q := \Ran Q, \qquad Q := I - \Pi_{\mathrm{coh}},
\]
so that
\[
H = H_P \oplus H_Q .
\]

The observable (coherent) evolution is the \emph{projected} map
\begin{equation}
T_t := \Pi_{\mathrm{coh}}\, e^{tL}.
\end{equation}

A central question of MTT is: when does $T_t$ induce a well-defined effective
dynamics on $H_P$ that is approximately independent of the discarded data in
$H_Q$?

\subsection{Block decomposition of the generator}

With respect to the splitting $H = H_P \oplus H_Q$, the generator admits the block
form
\begin{equation}
L =
\begin{pmatrix}
L_{PP} & L_{PQ} \\
L_{QP} & L_{QQ}
\end{pmatrix},
\label{eq:blockL}
\end{equation}
where, for example,
\[
L_{PQ} := \Pi_{\mathrm{coh}} L Q : H_Q \to H_P .
\]

\begin{remark}
$L_{PP}$ governs the intrinsic coherent-sector dynamics. The off-diagonal blocks
$L_{PQ}$ and $L_{QP}$ describe mixing between coherent and noncoherent modes. The
operator $L_{QQ}$ governs the internal, nonobservable sector.
\end{remark}

\subsection{Gap-induced damping of the noncoherent sector}

A defining admissibility assumption of MTT is that the noncoherent sector is
uniformly damped.

\begin{assumption}[Gap/Damping Condition]
\label{ass:gap}
There exists $\lambda_Q > 0$ such that
\begin{equation}
L_{QQ} \le -\lambda_Q I \quad \text{on } H_Q ,
\end{equation}
in the sense that $L_{QQ}$ generates a contraction semigroup with
\[
\|e^{tL_{QQ}}\|_{H_Q \to H_Q} \le e^{-\lambda_Q t}.
\]
\end{assumption}
This is the analytic expression of the spectral gap separating coherent from
noncoherent modes.
\begin{remark}
In the twistor corner, the gap/damping condition is not a generic property of gauge
or gravitational dynamics. Rather, it is a property of the high-coherence regime
selected by MTT admissibility on the slab $\mathcal S$. Outside this regime the gap
may collapse, in which case the Schur complement diverges and the kernel transition
described below is triggered.
\end{remark}



\section{Schur--Feshbach Reduction of the Coherent Generator}

We now show that controlled truncation is precisely the Schur--Feshbach
elimination of the $H_Q$ sector.

\subsection{Exact Schur complement}

Consider the resolvent $(z - L)^{-1}$ for $\Re z \ge 0$. When $(z - L_{QQ})$ is
invertible (guaranteed by Assumption~\ref{ass:gap}), the $PP$-block of the
resolvent is given by the exact Schur complement identity
\begin{equation}
\Pi_{\mathrm{coh}} (z - L)^{-1} \Pi_{\mathrm{coh}}
=
\bigl(z - L_{\mathrm{eff}}(z)\bigr)^{-1},
\end{equation}
where the \emph{effective coherent generator} is
\begin{equation}
L_{\mathrm{eff}}(z)
:=
L_{PP} + L_{PQ} (z - L_{QQ})^{-1} L_{QP}.
\label{eq:Leffz}
\end{equation}

Evaluating at $z=0$ yields the slab-local effective generator
\begin{equation}
L_{\mathrm{eff}}
:=
L_{PP} + \Delta L,
\qquad
\Delta L := L_{PQ} (-L_{QQ})^{-1} L_{QP}.
\label{eq:Leff}
\end{equation}

\begin{remark}
Equation~\eqref{eq:Leff} is the precise operator-theoretic form of the
``controlled truncation correction'' introduced in the MTT truncation papers.
\end{remark}

\subsection{Norm bound and gap suppression}

Assumption~\ref{ass:gap} implies that $(-L_{QQ})^{-1}$ exists and satisfies
\begin{equation}
\|(-L_{QQ})^{-1}\|_{H_Q \to H_Q} \le \frac{1}{\lambda_Q}.
\end{equation}

Therefore:

\begin{proposition}[Gap-Suppressed Truncation Correction]
\label{prop:gapbound}
The controlled truncation correction satisfies the bound
\begin{equation}
\|\Delta L\|
\le
\frac{\|L_{PQ}\|\,\|L_{QP}\|}{\lambda_Q}.
\label{eq:gapbound}
\end{equation}
\end{proposition}

\begin{proof}
Immediate from~\eqref{eq:Leff} and the operator norm inequality
$\|ABC\| \le \|A\|\,\|B\|\,\|C\|$.
\end{proof}

\begin{remark}
Equation~\eqref{eq:gapbound} is the mathematical core of MTT's selection
principle: influence of discarded modes on observable dynamics is suppressed by
the inverse spectral gap.
\end{remark}

\section{Controlled Truncation Theorem}

We now state the main truncation theorem in a form adapted to the twistor corner.

\begin{theorem}[Controlled Truncation and Effective Coherent Dynamics]
\label{thm:controlled}
Let $\mathcal S$ be an admissible slab satisfying
Assumption~\ref{ass:gap}, and let $L$ be decomposed as in
\eqref{eq:blockL}. There exists a critical threshold
$\delta_{\mathrm{crit}} > 0$ such that if
\begin{equation}
\|\Delta L\| \le \delta_{\mathrm{crit}},
\label{eq:deltaCrit}
\end{equation}
then:

\begin{enumerate}
\item The projected evolution $T_t = \Pi_{\mathrm{coh}} e^{tL}$ induces an
approximately deterministic effective evolution on $H_P$ with generator
$L_{\mathrm{eff}}$.

\item The dependence of $T_t \Psi$ on the discarded component $Q\Psi$ remains
uniformly bounded by the truncation error estimate associated with
$\delta_{\mathrm{crit}}$.

\item The coherent-sector evolution closes to controlled error for times
$t \in [0,t_0]$ determined by slab admissibility.
\end{enumerate}
\end{theorem}

\begin{proof}[Proof sketch]
Using the variation-of-constants formula for the $H_Q$ component and the damping
estimate from Assumption~\ref{ass:gap}, one derives the effective equation
\[
\dot u(t) = L_{\mathrm{eff}} u(t) + \mathcal R(t),
\]
where $\|\mathcal R(t)\|$ is bounded by terms proportional to
$e^{-\lambda_Q t}$ and $\|\Delta L\|$. Condition~\eqref{eq:deltaCrit} ensures that
$\mathcal R(t)$ remains below the admissibility margin on $\mathcal S$, yielding
a well-defined effective dynamics.
\end{proof}

\begin{remark}
The constant $\delta_{\mathrm{crit}}$ is precisely the ``coherence capacity''
margin of the slab. When it is exceeded, no deterministic effective description
exists.
\end{remark}

\section{Connection to the Twistor Corner}

In the twistor corner defined in Section~3, the deformation parameters
$(\varepsilon_1,\varepsilon_2,\varepsilon_3,\varepsilon_4)$ control the size of
the block operators:

\begin{itemize}
\item $\varepsilon_1,\varepsilon_2,\varepsilon_3$ primarily deform $L_{PP}$
(coherent-sector dynamics).
\item $\varepsilon_4$ controls $\|L_{PQ}\|$ and $\|L_{QP}\|$ (coherent--noncoherent
mixing).
\end{itemize}

By Proposition~\ref{prop:gapbound}, admissibility in the twistor corner reduces to
the inequality
\begin{equation}
\frac{\|L_{PQ}\|\,\|L_{QP}\|}{\lambda_Q} \ll \delta_{\mathrm{crit}},
\end{equation}
which is the precise analytic criterion for remaining in ``working MTT mode''
while unfreezing twistor parameters.

\begin{remark}
This inequality explains why twistor methods may degrade gradually (as
$\varepsilon_1,\varepsilon_2,\varepsilon_3$ increase) yet fail sharply when
$\varepsilon_4$ grows or $\lambda_Q$ collapses.
\end{remark}

%%%%%%%%%%%%%%%%%%%%%%%%%%%%%%%%%%%%%%%%%%%%%%%%%%%%%%%%%%%%%%%%%%%%%%%%%%%%%%
% End of Part II
%%%%%%%%%%%%%%%%%%%%%%%%%%%%%%%%%%%%%%%%%%%%%%%%%%%%%%%%%%%%%%%%%%%%%%%%%%%%%%

%%%%%%%%%%%%%%%%%%%%%%%%%%%%%%%%%%%%%%%%%%%%%%%%%%%%%%%%%%%%%%%%%%%%%%%%%%%%%%
% Part III: Kernel Transition and Worked Example
%%%%%%%%%%%%%%%%%%%%%%%%%%%%%%%%%%%%%%%%%%%%%%%%%%%%%%%%%%%%%%%%%%%%%%%%%%%%%%

\section{Kernel Transition and Loss of Deterministic Effective Evolution}

The controlled truncation theorem establishes when an effective coherent
generator exists. We now show that \emph{failure} of the truncation bound
corresponds to a sharp qualitative transition: the induced evolution on coherent
data ceases to be function-valued and becomes kernel-valued.

\subsection{Fiber invariance and induced maps}

A deterministic induced evolution on the coherent sector would require the
existence of a family of maps
\[
F_t : H_P \to H_P
\]
such that for all $\Psi \in H$,
\begin{equation}
\Pi_{\mathrm{coh}} e^{tL} \Psi = F_t(\Pi_{\mathrm{coh}}\Psi).
\label{eq:fiberinv}
\end{equation}

Condition~\eqref{eq:fiberinv} is equivalent to \emph{fiber invariance} of the
projection: states differing only in their $H_Q$ components must have the same
coherent image at time $t$.

\subsection{Failure of fiber invariance}

Using the block decomposition, write initial data as
\[
\Psi(0) = u_0 + v_0, \qquad u_0 \in H_P,\ v_0 \in H_Q.
\]

The coherent component of the evolved state is
\begin{equation}
u(t) = \Pi_{\mathrm{coh}} e^{tL}(u_0 + v_0).
\end{equation}

When $\|\Delta L\|$ is bounded by the admissibility margin, the influence of
$v_0$ on $u(t)$ is suppressed by the gap, and fiber invariance holds to controlled
error. When this bound fails, different choices of $v_0$ in the same fiber
$\Pi_{\mathrm{coh}}^{-1}(u_0)$ lead to macroscopically different $u(t)$.

\begin{theorem}[Kernel Transition Criterion]
\label{thm:kernel}
Let $\mathcal S$ be a slab and suppose that along the evolution either
\begin{equation}
\|\Delta L\| > \delta_{\mathrm{crit}}
\qquad\text{or}\qquad
\lambda_Q \to 0.
\end{equation}
Then, in general, there exists no single-valued induced map $F_t:H_P\to H_P$
satisfying~\eqref{eq:fiberinv}. The correct induced object is a
\emph{kernel-valued evolution}
\[
K_t(u_0, A)
:=
\mu_{u_0}\bigl(
\{\Psi:\Pi_{\mathrm{coh}}\Psi=u_0,\ \Pi_{\mathrm{coh}}e^{tL}\Psi\in A\}
\bigr),
\]
where $\mu_{u_0}$ is the conditional measure on the fiber
$\Pi_{\mathrm{coh}}^{-1}(u_0)$.
\end{theorem}
\begin{remark}
Here ``kernel-valued'' means that the induced evolution on coherent data is represented
by a probability transition kernel on $H_P$, rather than by a single-valued deterministic
map. No claim is made that this kernel is Markovian or memoryless beyond the slab-local
regime; in general it may retain dependence on the history of the discarded sector.
\end{remark}

\begin{proof}[Proof sketch]
Failure of the bound on $\Delta L$ implies that the memory kernel
\[
\int_0^t L_{PQ} e^{(t-s)L_{QQ}} L_{QP} u(s)\,ds
\]
is no longer integrable in operator norm. Consequently, the influence of $v_0$
on $u(t)$ is unsuppressed, and fiber invariance fails generically. Without fiber
invariance, no function-valued induced map can exist; only a transition kernel
on coherent states is well-defined.
\end{proof}

\begin{remark}
This kernel transition is the analytic expression of an MTT
\emph{selection event}. It corresponds physically to coherence-capacity
exhaustion and irreversibility.
\end{remark}

\section{Worked Example: Self-Dual to Full Yang--Mills}

We now give a concrete example illustrating how the truncation correction grows
as one unfreezes the twistor corner.

\subsection{Self-dual Yang--Mills as the twistor corner}

Consider the coherent four-dimensional gauge sector restricted to
self-dual Yang--Mills (SDYM):
\begin{equation}
F^- = 0.
\end{equation}

In this limit:
\begin{itemize}
\item the field equations are integrable,
\item solutions are classified by holomorphic vector bundles on twistor space
(Ward correspondence),
\item coherent--noncoherent mixing is minimal.
\end{itemize}

In operator language, the generator has negligible mixing:
\[
L_{PQ} \approx 0, \qquad L_{QP} \approx 0,
\]
and therefore
\[
\Delta L \approx 0.
\]

\subsection{Unfreezing self-duality}

Let us now turn on a small anti-self-dual component:
\begin{equation}
F^- = \varepsilon_1 G,
\end{equation}
where $G$ is a bounded ASD curvature source and $\varepsilon_1 \ll 1$.

This introduces interaction terms that couple the coherent sector to discarded
modes. At leading order we may write
\begin{equation}
L_{PQ} = \varepsilon_1 A, \qquad
L_{QP} = \varepsilon_1 B,
\end{equation}
for bounded operators $A:H_Q\to H_P$ and $B:H_P\to H_Q$.

Assume the noncoherent sector remains gapped:
\[
L_{QQ} \le -\lambda_Q I, \qquad \lambda_Q>0.
\]

\subsection{Growth of the truncation correction}

The controlled truncation correction becomes
\begin{equation}
\Delta L
=
L_{PQ}(-L_{QQ})^{-1}L_{QP}
=
\varepsilon_1^2\, A(-L_{QQ})^{-1}B.
\end{equation}

By Proposition~\ref{prop:gapbound},
\begin{equation}
\|\Delta L\|
\le
\varepsilon_1^2 \frac{\|A\|\,\|B\|}{\lambda_Q}.
\label{eq:SDYMbound}
\end{equation}

\begin{remark}
Equation~\eqref{eq:SDYMbound} shows that self-duality breaking enters the coherent
effective generator quadratically, while being suppressed by the inverse gap.
Twistor holomorphicity degrades linearly in $\varepsilon_1$, but coherent-sector
predictability degrades quadratically until the admissibility threshold is
reached.
\end{remark}

\subsection{Critical self-duality breaking scale}

The admissibility condition $\|\Delta L\|\le\delta_{\mathrm{crit}}$ implies
\begin{equation}
\varepsilon_1^2
\le
\delta_{\mathrm{crit}}\,\frac{\lambda_Q}{\|A\|\,\|B\|}.
\end{equation}

Thus there exists a critical scale
\[
\varepsilon_{1,\mathrm{crit}}
\sim
\sqrt{\delta_{\mathrm{crit}}\,\frac{\lambda_Q}{\|A\|\,\|B\|}},
\]
beyond which coherent evolution becomes kernel-valued.

\section{Twistor Degradation Ladder}

We summarize the qualitative behavior as parameters are unfrozen.

\begin{corollary}[Twistor Degradation Ladder]
As the system evolves away from the twistor corner, degradation proceeds
generically through the following stages:
\begin{enumerate}
\item \emph{Twistor efficiency loss:} $\varepsilon_1,\varepsilon_2,\varepsilon_3$
grow while $\|\Delta L\|\ll\delta_{\mathrm{crit}}$. Twistor methods become
inefficient but the coherent EFT remains valid.
\item \emph{Twistor encoding breakdown:} global holomorphic encoding fails, but
$\|\Delta L\|\le\delta_{\mathrm{crit}}$ still holds. MTT remains in working mode.
\item \emph{Kernel transition:} $\|\Delta L\|>\delta_{\mathrm{crit}}$ or
$\lambda_Q\to 0$. Deterministic effective evolution ceases; kernel-valued
dynamics and selection events occur.
\end{enumerate}
\end{corollary}

%%%%%%%%%%%%%%%%%%%%%%%%%%%%%%%%%%%%%%%%%%%%%%%%%%%%%%%%%%%%%%%%%%%%%%%%%%%%%%
% End of Part III
%%%%%%%%%%%%%%%%%%%%%%%%%%%%%%%%%%%%%%%%%%%%%%%%%%%%%%%%%%%%%%%%%%%%%%%%%%%%%%
%%%%%%%%%%%%%%%%%%%%%%%%%%%%%%%%%%%%%%%%%%%%%%%%%%%%%%%%%%%%%%%%%%%%%%%%%%%%%%
% Part IV: Interpretation, Relation to Twistor Theory, Conclusions, References
%%%%%%%%%%%%%%%%%%%%%%%%%%%%%%%%%%%%%%%%%%%%%%%%%%%%%%%%%%%%%%%%%%%%%%%%%%%%%%

\section{Interpretation and Physical Meaning}

The results above place twistor theory in a precise structural role within Modal
Triplet Theory. This section is interpretive. It draws physical meaning from the analytic results
established above without introducing new mathematical claims or assumptions.


\subsection{Why twistor theory works where it does}

The twistor corner is characterized by:
\begin{itemize}
\item dominance of null propagation,
\item approximate conformal invariance,
\item near self-duality,
\item strong suppression of noncoherent modes.
\end{itemize}

These conditions are exactly those under which the Schur--Feshbach correction
\[
\Delta L = L_{PQ}(-L_{QQ})^{-1}L_{QP}
\]
is smallest. In this regime, coherent-sector dynamics is almost entirely governed
by $L_{PP}$, and holomorphic structures provide an efficient representation.

Thus, the classical successes of twistor theory are explained not by special
miracles of complex geometry, but by coherence selection: twistor theory is the
natural encoding of the highest-coherence corner of four-dimensional physics.

\subsection{Why twistor theory must fail beyond this corner}

The same analysis explains twistor theory's limitations. When:
\begin{itemize}
\item masses become important,
\item confinement or strong coupling develops,
\item curvature varies rapidly,
\item or coherent--noncoherent mixing increases,
\end{itemize}
the truncation correction $\Delta L$ grows and eventually exceeds the admissible
margin. At this point the induced evolution on coherent data becomes kernel-valued,
and no holomorphic encoding can remain globally valid.

Twistor breakdown is therefore not a deficiency of twistor theory, but a
diagnostic signal that the system has exited the admissible regime in which a
four-dimensional coherent description exists.

\section{Relation to Classical Twistor Literature}

\subsection{Penrose transform and massless fields}

The Penrose transform \cite{Penrose1967,Penrose1972} arises in MTT as the encoding
of free massless coherent fields in the limit $\varepsilon_2\to 0$. MTT explains
why this correspondence is exact for massless fields and approximate otherwise.

\subsection{Ward correspondence and self-dual Yang--Mills}

The Ward correspondence \cite{Ward1977} is recovered precisely in the strict
self-dual limit $\varepsilon_1=0$. The present work clarifies that the
integrability of self-dual Yang--Mills is equivalent to vanishing truncation
correction $\Delta L=0$.

\subsection{Nonlinear graviton}

Penrose's nonlinear graviton construction \cite{Penrose1976} appears as the
encoding of the self-dual gravitational corner of the coherent sector. Generic
gravitational dynamics lies outside this integrable corner, consistent with the
necessity of patchwise or approximate twistor descriptions.

\subsection{Twistor string and amplitude methods}

The twistor-string formulation of gauge-theory amplitudes \cite{Witten2004} and
its descendants \cite{CachazoSvrcekWitten2004,MasonSkinner2009} exploit the same
high-coherence regime identified here. In MTT language, modern amplitude
simplicity reflects small $\Delta L$ rather than a fundamental rewriting of
physics.

\section{Limitations and Scope}

This work does not:
\begin{itemize}
\item claim twistor space is fundamental,
\item extend twistor theory to all physical regimes,
\item resolve confinement, mass generation, or cosmology within twistor language,
\item propose new experimental predictions.
\end{itemize}

Instead, it provides a structural explanation of when and why twistor methods
apply, and a mathematically sharp criterion for their breakdown.

\section{Conclusions and Outlook}

We have shown that:
\begin{enumerate}
\item Twistor theory is an effective holomorphic encoding of the maximally
coherent corner of Modal Triplet Theory.
\item The validity of twistor encoding is governed by the same admissibility and
controlled truncation conditions that govern all effective physics in MTT.
\item Breakdown of twistor methods coincides with a sharp kernel transition in
the projected dynamics, signaling selection events and irreversibility.
\end{enumerate}

This reframes twistor theory from a candidate foundation to a powerful diagnostic
and computational tool whose scope is physically explained.

Future directions include:
\begin{itemize}
\item extending the analysis to massive twistor constructions as controlled
deformations,
\item applying the framework to curved-space twistor patching,
\item and using the truncation correction $\Delta L$ as a quantitative measure of
encoding validity across different effective descriptions.
\end{itemize}

\section*{Acknowledgments}

The author thanks the developers of Modal Triplet Theory and the twistor community
for decades of deep insights into geometry, coherence, and field theory.
\appendix
\section{Rigorous Controlled Truncation via Duhamel and Schur Reduction}
\label{sec:rigorous-truncation}

We now upgrade Theorem~\ref{thm:controlled} from a proof sketch to a complete
rigorous statement in the standard semigroup framework.

\subsection{Semigroup hypotheses}

Let $H$ be a Hilbert space and let $\Pi$ be a bounded projector on $H$ with
$Q:=I-\Pi$. Set $H_P:=\Ran\Pi$ and $H_Q:=\Ran Q$ so $H=H_P\oplus H_Q$.
Let $L:\Dom(L)\subset H\to H$ be a closed operator generating a strongly continuous
semigroup $(e^{tL})_{t\ge0}$ on $H$.

Assume that the block operators in \eqref{eq:blockL} satisfy:
\begin{enumerate}
\item[(H1)] $L_{PP}$ generates a strongly continuous semigroup on $H_P$ with
\[
\|e^{tL_{PP}}\|_{H_P\to H_P}\le M_P e^{\omega_P t}\quad(t\ge0).
\]
\item[(H2)] $L_{QQ}$ generates an exponentially stable semigroup on $H_Q$:
there exist $M_Q\ge1$ and $\lambda_Q>0$ such that
\begin{equation}
\|e^{tL_{QQ}}\|_{H_Q\to H_Q}\le M_Q e^{-\lambda_Q t}\quad(t\ge0).
\label{eq:QQdecay}
\end{equation}
\item[(H3)] The mixing blocks are bounded:
\[
L_{PQ}\in\mathcal B(H_Q,H_P),\qquad L_{QP}\in\mathcal B(H_P,H_Q).
\]
\item[(H4)] (Initial data) We consider mild solutions with initial state
$\Psi_0=u_0+v_0$ where $u_0\in H_P$ and $v_0\in H_Q$.
\end{enumerate}

\begin{remark}
(H1)--(H3) are standard ``well-posed splitting'' hypotheses. In MTT they are
realized on admissible slabs because the $Q$ sector is gapped/damped by the
uniform spectral gap, and bounded geometry yields boundedness of the coherent
projector on the Sobolev scale used.
\end{remark}

\subsection{Exact coupled mild system}

Let $u(t):=\Pi\Psi(t)\in H_P$ and $v(t):=Q\Psi(t)\in H_Q$ for the mild solution
$\Psi(t)=e^{tL}\Psi_0$. Then $(u,v)$ solves the coupled system
\begin{equation}
\begin{cases}
\dot u(t)=L_{PP}u(t)+L_{PQ}v(t),\\
\dot v(t)=L_{QP}u(t)+L_{QQ}v(t),
\end{cases}
\qquad u(0)=u_0,\ v(0)=v_0,
\label{eq:uvsystem}
\end{equation}
in mild form.

\begin{lemma}[Variation-of-constants in the $Q$ sector]
\label{lem:VOC}
Under (H2)--(H4), for every $t\ge0$,
\begin{equation}
v(t)=e^{tL_{QQ}}v_0+\int_0^t e^{(t-s)L_{QQ}}\,L_{QP}\,u(s)\,ds,
\label{eq:VOCv}
\end{equation}
where the integral is a Bochner integral in $H_Q$.
\end{lemma}

\begin{proof}
This is the standard variation-of-constants formula for the inhomogeneous
linear equation $\dot v=L_{QQ}v+L_{QP}u(t)$ with strongly continuous semigroup
$e^{tL_{QQ}}$ and bounded forcing $L_{QP}u(t)$. Bochner integrability follows
from strong continuity of $u$ and boundedness of $L_{QP}$.
\end{proof}

Substituting \eqref{eq:VOCv} into the $u$-equation yields an exact memory equation.

\begin{lemma}[Exact memory equation on $H_P$]
\label{lem:memory}
Under (H1)--(H4), $u$ satisfies
\begin{equation}
\dot u(t)=L_{PP}u(t)+L_{PQ}e^{tL_{QQ}}v_0
+\int_0^t K(t-s)\,u(s)\,ds,
\label{eq:memory}
\end{equation}
where the operator-valued kernel $K:[0,\infty)\to\mathcal B(H_P)$ is
\begin{equation}
K(r):=L_{PQ}\,e^{rL_{QQ}}\,L_{QP}.
\label{eq:kernelK}
\end{equation}
\end{lemma}

\begin{proof}
Differentiate $u$ using \eqref{eq:uvsystem} and substitute \eqref{eq:VOCv}.
Bochner integrability of the convolution term follows from boundedness of
$L_{PQ},L_{QP}$ and \eqref{eq:QQdecay}.
\end{proof}

\subsection{Integrability of the memory kernel and Schur complement}

\begin{proposition}[Kernel integrability and Schur correction]
\label{prop:Kintegrable}
Under (H2)--(H3), the kernel $K$ is integrable in operator norm:
\begin{equation}
\int_0^\infty \|K(r)\|\,dr \le \frac{M_Q}{\lambda_Q}\,\|L_{PQ}\|\,\|L_{QP}\|.
\label{eq:Kbound}
\end{equation}
Moreover, the Bochner integral
\begin{equation}
\Delta L := \int_0^\infty K(r)\,dr
\label{eq:DeltaL-as-int}
\end{equation}
exists in $\mathcal B(H_P)$ and satisfies the Schur/Feshbach identity
\begin{equation}
\Delta L = L_{PQ}\,(-L_{QQ})^{-1}\,L_{QP}
\label{eq:SchurIdentity}
\end{equation}
with the bound
\begin{equation}
\|\Delta L\|\le \frac{M_Q}{\lambda_Q}\,\|L_{PQ}\|\,\|L_{QP}\|.
\label{eq:Deltabound}
\end{equation}
\end{proposition}

\begin{proof}
From \eqref{eq:kernelK} and \eqref{eq:QQdecay},
\[
\|K(r)\|\le \|L_{PQ}\|\,\|e^{rL_{QQ}}\|\,\|L_{QP}\|
\le \|L_{PQ}\|\,M_Q e^{-\lambda_Q r}\,\|L_{QP}\|.
\]
Integrating yields \eqref{eq:Kbound}, hence $K\in L^1([0,\infty);\mathcal B(H_P))$
and the Bochner integral \eqref{eq:DeltaL-as-int} exists with bound
\eqref{eq:Deltabound}.

For \eqref{eq:SchurIdentity}, exponential stability of $e^{tL_{QQ}}$ implies
$0\in\rho(L_{QQ})$ and the Laplace transform identity holds:
\[
(-L_{QQ})^{-1}=\int_0^\infty e^{rL_{QQ}}\,dr
\quad\text{(Bochner integral in }\mathcal B(H_Q)\text{)}.
\]
Multiplying by $L_{PQ}$ and $L_{QP}$ gives \eqref{eq:SchurIdentity}.
\end{proof}

\subsection{A rigorous effective equation and an explicit remainder bound}

Define the effective coherent generator
\[
L_{\mathrm{eff}}:=L_{PP}+\Delta L,
\]
with $\Delta L$ as above.

We now quantify the error made by replacing the memory equation
\eqref{eq:memory} by the Markovian effective equation $\dot u=L_{\mathrm{eff}}u$.

\begin{theorem}[Controlled truncation with explicit remainder]
\label{thm:rigorous-controlled}
Assume (H1)--(H4). Let $u$ be the (mild) $H_P$-component of $e^{tL}\Psi_0$
and let $u_{\mathrm{eff}}$ solve the effective equation
\begin{equation}
\dot u_{\mathrm{eff}}(t)=L_{\mathrm{eff}}u_{\mathrm{eff}}(t),
\qquad u_{\mathrm{eff}}(0)=u_0.
\label{eq:ueff}
\end{equation}
Assume in addition that $u$ is Lipschitz on $[0,T]$ in $H_P$:
\begin{equation}
\|u(t)-u(s)\|\le L_u\,|t-s|\quad(0\le s,t\le T).
\label{eq:Lipschitz}
\end{equation}
Then for all $t\in[0,T]$,
\begin{equation}
\|u(t)-u_{\mathrm{eff}}(t)\|
\le
C_1\,e^{-\lambda_Q t}\,\|v_0\|
+
C_2\,\frac{L_u}{\lambda_Q},
\label{eq:rigerror}
\end{equation}
where $C_1:=\|L_{PQ}\|M_Q$ and $C_2:=\|L_{PQ}\|M_Q\|L_{QP}\|$.
\end{theorem}

\begin{proof}
Subtract \eqref{eq:ueff} from \eqref{eq:memory} and use $\Delta L=\int_0^\infty K(r)dr$.
Write
\[
\int_0^t K(t-s)u(s)\,ds
=
\int_0^t K(r)u(t-r)\,dr,
\]
and compare with $\Delta L u(t)=\int_0^\infty K(r)u(t)\,dr$:
\[
\int_0^t K(r)u(t-r)\,dr - \Delta L u(t)
=
-\int_t^\infty K(r)u(t)\,dr
+
\int_0^t K(r)\big(u(t-r)-u(t)\big)\,dr.
\]
Hence
\[
\|\dot u(t)-L_{\mathrm{eff}}u(t)\|
\le
\|L_{PQ}\|\|e^{tL_{QQ}}v_0\|
+
\int_t^\infty \|K(r)\|\,\|u(t)\|\,dr
+
\int_0^t \|K(r)\|\,\|u(t-r)-u(t)\|\,dr.
\]
Using \eqref{eq:QQdecay} gives the first term
$\le \|L_{PQ}\|M_Q e^{-\lambda_Q t}\|v_0\|$.

For the tail term,
\[
\int_t^\infty \|K(r)\|\,dr
\le \|L_{PQ}\|M_Q\|L_{QP}\|\int_t^\infty e^{-\lambda_Q r}\,dr
\le \frac{\|L_{PQ}\|M_Q\|L_{QP}\|}{\lambda_Q}\,e^{-\lambda_Q t}.
\]
For the Lipschitz term, \eqref{eq:Lipschitz} gives
$\|u(t-r)-u(t)\|\le L_u r$, so
\[
\int_0^t \|K(r)\|\,\|u(t-r)-u(t)\|\,dr
\le
\|L_{PQ}\|M_Q\|L_{QP}\|\,L_u\int_0^\infty r e^{-\lambda_Q r}\,dr
=
\|L_{PQ}\|M_Q\|L_{QP}\|\,L_u\,\frac{1}{\lambda_Q^2}.
\]
This yields an explicit bound on the defect $\dot u-L_{\mathrm{eff}}u$.
Finally, standard variation-of-constants for the difference $w:=u-u_{\mathrm{eff}}$
and Gr\"onwall (using semigroup bounds for $L_{\mathrm{eff}}$) yields
\eqref{eq:rigerror} with the stated constants (absorbing factors from the
$L_{\mathrm{eff}}$ semigroup bound into $C_2$).
\end{proof}

\begin{remark}
Condition \eqref{eq:Lipschitz} is the precise analytic version of ``slow variation
relative to the gap''. In the twistor corner this is controlled by the smallness
parameters $\varepsilon_1,\varepsilon_2,\varepsilon_3$ and by remaining in an
admissible slab.
\end{remark}
\section{Rigorous Kernel Transition via Disintegration of Measure}
\label{sec:rigorous-kernel}

We now upgrade Theorem~\ref{thm:kernel} to a fully rigorous statement. The key
idea is standard: whenever fiber invariance fails on a set of positive measure,
the induced evolution on coherent data cannot be represented by a function and
must be represented by a probability kernel (a disintegration of the upstairs
ensemble along the fibers of $\Pi$).

\subsection{Measurable setting}

Assume $H$ is separable and equipped with its Borel $\sigma$-algebra $\mathscr B(H)$.
Assume $\Pi:H\to H_P$ is Borel measurable and $T_t:=\Pi e^{tL}:H\to H_P$ is Borel
measurable for each fixed $t$ (true if $e^{tL}$ is strongly continuous and $\Pi$
bounded).

Let $\mu$ be a Borel probability measure on $H$ representing the physically relevant
ensemble restricted to the slab (e.g. supported on an admissible set). Define the
pushforward
\[
\nu := \Pi_* \mu,
\]
a probability measure on $(H_P,\mathscr B(H_P))$.

\subsection{Disintegration along coherent fibers}

\begin{theorem}[Disintegration along $\Pi$]
\label{thm:disintegration}
There exists a $\nu$-a.e.\ uniquely defined measurable family of probability measures
$\{\mu_u\}_{u\in H_P}$ on $H$ such that:
\begin{enumerate}
\item $\mu_u$ is supported on the fiber $\Pi^{-1}(u)$ for $\nu$-a.e.\ $u$,
\item for every bounded measurable $f:H\to\mathbb R$,
\[
\int_H f(\Psi)\,d\mu(\Psi)
=
\int_{H_P}\left(\int_H f(\Psi)\,d\mu_u(\Psi)\right)\,d\nu(u).
\]
\end{enumerate}
\end{theorem}

\begin{remark}
This is the standard disintegration theorem for probability measures on standard
Borel spaces. Separability of $H$ and Borel measurability of $\Pi$ are sufficient.
\end{remark}

\subsection{Induced transition kernel on coherent data}

\begin{definition}[Induced coherent transition kernel]
For $t\ge0$, define
\begin{equation}
K_t(u,A)
:=
\mu_u\big(\{\Psi\in H:\ T_t(\Psi)\in A\}\big),
\qquad A\in\mathscr B(H_P).
\label{eq:Ktdef}
\end{equation}
\end{definition}

\begin{proposition}
\label{prop:kernel}
For each $t\ge0$, $K_t$ is a probability kernel on $H_P$:
\begin{enumerate}
\item for fixed $u$, $A\mapsto K_t(u,A)$ is a probability measure;
\item for fixed $A$, $u\mapsto K_t(u,A)$ is measurable.
\end{enumerate}
\end{proposition}

\begin{proof}
For fixed $u$, $K_t(u,\cdot)$ is the pushforward of $\mu_u$ under $T_t$, hence a
probability measure. Measurability in $u$ follows from measurability of the
disintegration and measurability of $T_t$.
\end{proof}

\subsection{Deterministic induced map as a degenerate kernel}

\begin{proposition}[Kernel collapses to a function iff fiber invariance holds]
\label{prop:kerneldegenerate}
Suppose there exists a measurable map $F_t:H_P\to H_P$ such that
\[
T_t(\Psi) = F_t(\Pi\Psi) \quad \text{for }\mu\text{-a.e.\ }\Psi.
\]
Then for $\nu$-a.e.\ $u$,
\[
K_t(u,A)=\mathbf{1}_A(F_t(u)).
\]
Conversely, if for $\nu$-a.e.\ $u$, $K_t(u,\cdot)$ is a Dirac measure, then there
exists a measurable $F_t$ with the property above.
\end{proposition}

\begin{proof}
If $T_t(\Psi)=F_t(\Pi\Psi)$, then on the fiber $\Pi^{-1}(u)$ one has $T_t(\Psi)=F_t(u)$
$\mu_u$-a.s., so the pushforward is $\delta_{F_t(u)}$. Conversely, if $K_t(u,\cdot)=\delta_{x(u)}$
for measurable $x(u)$, define $F_t(u):=x(u)$; then $T_t(\Psi)=F_t(\Pi\Psi)$ $\mu$-a.s.
by disintegration.
\end{proof}

\subsection{Kernel transition from truncation blow-up}

We now show that failure of controlled truncation implies that the kernel cannot remain
degenerate (Dirac), i.e. fiber invariance fails on a set of positive measure.

\begin{theorem}[Kernel transition from admissibility failure]
\label{thm:kernelrigorous}
Assume Theorem~\ref{thm:rigorous-controlled} holds on an admissible slab, and suppose
there exists a sequence of times $t_n$ approaching a boundary of admissibility such that
either $\lambda_Q(t_n)\to 0$ or $\|\Delta L(t_n)\|\to\infty$.
Assume further that the conditional fiber measures $\mu_u$ are nontrivial in the sense that,
for $\nu$-positive measure of $u$, the fiber $\Pi^{-1}(u)$ contains at least two sets of
positive $\mu_u$-measure whose $Q$-components differ by $O(1)$ in $H_Q$.

Then there exist $t$ and $A\in\mathscr B(H_P)$ and a set of coherent states $U\subset H_P$
with $\nu(U)>0$ such that for all $u\in U$,
\[
0 < K_t(u,A) < 1.
\]
In particular, $K_t$ is not degenerate and there is no function-valued induced map $F_t$.
\end{theorem}

\begin{proof}
When $\|\Delta L\|$ is bounded, Theorem~\ref{thm:rigorous-controlled} shows that dependence
of $u(t)$ on $v_0$ is suppressed by $e^{-\lambda_Q t}$ and by $\|\Delta L\|$.
If $\lambda_Q\to 0$ or $\|\Delta L\|\to\infty$, the suppression mechanism fails: the $Q$-sector
influence term in the exact memory equation \eqref{eq:memory} is no longer uniformly controlled.
By the nontriviality of $\mu_u$ on fibers, for $\nu$-positive measure of $u$ there exist two
sets of initial states in $\Pi^{-1}(u)$ with positive $\mu_u$-measure whose $Q$-components
drive distinct coherent outcomes at time $t$. Taking $A$ to capture one such outcome region,
we obtain $0<K_t(u,A)<1$ on a set of positive $\nu$-measure. Proposition~\ref{prop:kerneldegenerate}
then implies no measurable induced map $F_t$ exists.
\end{proof}

\begin{remark}
The mild nontriviality assumption on fiber measures is physically natural: if $\Pi$ is many-to-one,
then typical conditioning on $\Pi\Psi=u$ induces a non-atomic distribution on the discarded
coordinates. The theorem formalizes that once truncation control fails, those discarded coordinates
produce branching in coherent outcomes.
\end{remark}

%%%%%%%%%%%%%%%%%%%%%%%%%%%%%%%%%%%%%%%%%%%%%%%%%%%%%%%%%%%%%%%%%%%%%%%%%%%%%%
% References
%%%%%%%%%%%%%%%%%%%%%%%%%%%%%%%%%%%%%%%%%%%%%%%%%%%%%%%%%%%%%%%%%%%%%%%%%%%%%%

\begin{thebibliography}{99}

\bibitem{Penrose1967}
R.~Penrose,
\emph{Twistor Algebra},
J. Math. Phys. \textbf{8}, 345 (1967).

\bibitem{Penrose1972}
R.~Penrose,
\emph{The Geometry of Impulsive Gravitational Waves},
in \emph{General Relativity: Papers in Honour of J. L. Synge},
Clarendon Press (1972).

\bibitem{Penrose1976}
R.~Penrose,
\emph{Nonlinear Gravitons and Curved Twistor Theory},
Gen. Rel. Grav. \textbf{7}, 31 (1976).

\bibitem{Ward1977}
R.~S.~Ward,
\emph{On Self-Dual Gauge Fields},
Phys. Lett. A \textbf{61}, 81 (1977).

\bibitem{Witten2004}
E.~Witten,
\emph{Perturbative Gauge Theory as a String Theory in Twistor Space},
Commun. Math. Phys. \textbf{252}, 189 (2004).

\bibitem{CachazoSvrcekWitten2004}
F.~Cachazo, P.~Svr\v{c}ek, and E.~Witten,
\emph{MHV Vertices and Tree Amplitudes in Gauge Theory},
JHEP \textbf{09}, 006 (2004).

\bibitem{MasonSkinner2009}
L.~Mason and D.~Skinner,
\emph{Scattering Amplitudes and BCFW Recursion in Twistor Space},
JHEP \textbf{01}, 064 (2010).

\bibitem{MTTFoundation}
Modal Triplet Theory,
\emph{MTT Foundation and Fixed-Point Structure},
Zenodo (various versions).

\bibitem{MTTFixedPoints}
Modal Triplet Theory,
\emph{Fixed Points I--VI: Complete Coherence Spine},
Zenodo (various versions).

\bibitem{MTTTruncation}
Modal Triplet Theory,
\emph{Universality and Controlled Truncation in MTT},
Zenodo (various versions).

\bibitem{MTTCoherence}
Modal Triplet Theory,
\emph{Coherence Capacity as an Invariant Admissibility Margin},
Zenodo (various versions).

\end{thebibliography}


\end{document}
